% Generated draft: requires mathematical and source review.
\documentclass[11pt]{article}
\usepackage[margin=25mm]{geometry}
\usepackage{fontspec}
\setmainfont{Latin Modern Roman}
\usepackage{amsmath,amssymb,mathtools,unicode-math}
\setmathfont{Latin Modern Math}
\usepackage{enumitem,needspace,xurl,hyperref}
\setlength{\parindent}{0pt}
\setlength{\parskip}{5pt}
\setlength{\emergencystretch}{3em}
\newcommand{\N}{\mathbb{N}}
\newcommand{\Z}{\mathbb{Z}}
\newcommand{\Q}{\mathbb{Q}}
\newcommand{\R}{\mathbb{R}}
\newcommand{\C}{\mathbb{C}}
\newcommand{\F}{\mathbb{F}}
\newcommand{\A}{\mathbb{A}}
\newcommand{\PP}{\mathbb{P}}
\newcommand{\E}{\mathbb{E}}
\newcommand{\eps}{\varepsilon}
\newcommand{\dd}{\,\mathrm{d}}
\begin{document}
\section*{UnsolvedMath ID: 30001518}
\textbf{Source problem code:} catalogue entry
\subsection{Definitions and mathematical statement}
A billiard trajectory travels along straight lines between reflections. At a regular reflecting boundary point with unit normal \(n\), specular reflection changes an incoming unit velocity \(v\) to
\[
 v-2\langle v,n\rangle n.
\]
When a trajectory has incoming and outgoing asymptotic unit velocities \(v_{\mathrm{in}}\) and \(v_{\mathrm{out}}\), retroreflection means
\[
 v_{\mathrm{out}}=-v_{\mathrm{in}}.
\]
This condition concerns reversal of direction; it does not require the outgoing ray to retrace the incoming line.

The supplied question asks whether a \emph{perfect billiard retroreflector} exists: a body that reverses the direction of every incident particle or ray. Schematically, if \(\mathcal B\) denotes the intended class of bodies and \(\mathcal I(B)\) denotes all incident ray data for \(B\), the requested property is
\[
 \exists B\in\mathcal B\quad
 \forall z\in\mathcal I(B),\qquad
 v_{\mathrm{out}}(B,z)\text{ exists and }
 v_{\mathrm{out}}(B,z)=-v_{\mathrm{in}}(z).
\]
Here the existence of an outgoing velocity expresses the requirement that each incident ray actually emerge with reversed direction.

\emph{Unresolved qualifications.} The record does not specify the ambient dimension, the admissible class \(\mathcal B\), or a complete definition of \(\mathcal I(B)\). In particular, it does not state boundedness, boundary regularity, allowable cavities or reflecting components, or conventions for tangential impacts, singular boundary points, trapped trajectories, and infinitely many reflections. Thus the displayed formula records the intended universal property but is not a fully specified existence problem until those choices are supplied. The source says \emph{every} incident particle or ray; replacing this by almost every incident datum would change the stated quantifier and is not justified by this record.
\subsection{Short English statement}
Does there exist a billiard body that sends every incident particle or ray back with exactly the opposite velocity direction? The source leaves the dimension, allowable bodies, and exceptional-trajectory conventions unspecified, so these require clarification before the question has a fully precise formulation.
\subsection{Sources}
\begin{itemize}
\item Local UnsolvedMath registry: supplied problem record with permanent ID \(30001518\), source code \texttt{catalogue entry}, in the Oberwolfach Reports Open Problems collection. The record marks the statement as exact and internally readable; this is supplied metadata, not independent mathematical verification. No individual proposer is identified; the supplied proposed year is 2010.
\item Supplied report citation: \emph{Calculus of Variations (2011)}, Oberwolfach Reports, DOI 10.4171/owr/2010/31, pp.~1827--1884. \url{https://doi.org/10.4171/owr/2010/31}. The background identifies Oberwolfach Reports as documenting mathematical workshops and their open-problem discussions. The citation's 2011 label and its 2010 DOI/year metadata are retained as supplied.
\item Supplied literature source: Plakhov, \emph{Mathematical retroreflectors}. \url{https://arxiv.org/abs/1005.3466}.
\item Supplied literature source: \emph{Oberwolfach Report 31, Calculus of Variations}. \url{https://doi.org/10.4171/owr/2010/31}.
\item Literature assessment dated \texttt{2026-08-21}: the supplied curation labels the problem open and states that the report described perfect billiard retroreflectors as unknown, with only asymptotically perfect constructions available, and that the survey likewise presents no perfect construction or impossibility theorem. This is supplied dated curation derived from the OWR record, not an independently verified literature assessment; the time-sensitive status and cited claims require independent checking.
\end{itemize}
\end{document}
